\section{Introduction}

The imminent arrival of 6G networks promises to enable next-generation edge intelligence with ultra-low latency, high reliability, and context-aware services including extended reality, autonomous systems, and tactile internet applications~\cite{Lou2022_CostEffectiveScheduling, Wang2025_AverageReliabilityOptimal}. These use cases demand fine-grained orchestration of computing and data services across heterogeneous edge nodes. Traditional cloud-based approaches cannot meet such requirements due to latency bottlenecks, reinforcing the importance of edge computing to bring processing closer to end users~\cite{Xu2023_FedLC}.

The dynamic and resource-constrained nature of edge environments combined with energy constraints, privacy requirements, and workload variability introduces substantial orchestration challenges~\cite{Yuan2025}. Recent studies have investigated reinforcement learning for adaptive orchestration, allowing real-time learning of placement and migration strategies~\cite{Yuan2025}. However, centralized RL methods suffer from scalability limitations and cannot preserve user privacy. Furthermore, most RL-based orchestration approaches disregard semantic content in service tasks, leading to suboptimal placements in latency-sensitive or context-dependent scenarios such as extended reality or real-time analytics~\cite{Yu2025}.

Federated reinforcement learning provides a promising alternative by distributing policy learning across edge nodes without sharing raw data~\cite{Chen2025_FRL}. While FRL frameworks have shown promise in handling privacy and scalability, they frequently overlook two critical dimensions: semantic task awareness and network topology encoding. Modern edge workloads often include semantic information such as task intent or content category, which if properly modeled can significantly enhance placement fidelity and decision stability~\cite{Wu2024_FRLEdgeComputing}. At the same time, orchestration effectiveness depends on network structure, including link quality, proximity, and cluster connectivity. Yet most FRL methods treat edge nodes as unstructured agents without considering topological relationships.

Graph neural networks offer an effective means to encode network topology~\cite{Sagar2025_HAFedRL}, but their integration into federated, real-time orchestration remains limited. Moreover, frequent global synchronization in federated learning systems often results in high communication overhead and latency. Hierarchical FL frameworks reduce this burden by performing local aggregation within clusters before periodic global updates~\cite{Ma2024_HierarchicalFL}. However, few approaches combine hierarchical control with semantic embeddings and graph-based topology encoding in a unified orchestration framework.

Existing studies also suffer from limited experimental scope. Many evaluations use synthetic environments with minimal orchestration realism or baseline diversity. Large-scale scenarios are rarely validated beyond small-scale simulations, and systematic scalability analysis using mathematical extrapolation from empirically validated base experiments remains uncommon in the federated edge orchestration literature.

\noindent\textbf{Contributions.} We propose \textbf{FedSemGNN}, a federated reinforcement learning framework designed for semantic-aware, topology-informed, and hierarchical orchestration in 6G edge environments. The key contributions of this work include:

\begin{itemize}
	\item \textbf{Semantic task embedding with continual learning:} Service requests are embedded into 16-dimensional semantic spaces using a custom continual learning encoder that extracts features from resource demands (CPU, memory, network), service types, and temporal patterns. The encoder employs elastic weight consolidation with regularization parameter 0.4 to prevent catastrophic forgetting while adapting to emerging workload characteristics. Semantic matching uses cosine similarity to enable intent-aware placement beyond numerical resource matching.
	
	\item \textbf{Graph convolutional network topology encoding:} A two-layer GCN with symmetric normalization captures structural relationships including bandwidth, affinity, and proximity. The architecture processes node features through 128 to 64 to 64 dimensional transformations with ReLU activation, allowing placement decisions that consider network-wide constraints and multi-hop dependencies across mesh-connected edge infrastructure.
	
	\item \textbf{Hierarchical federated PPO with dual-level synchronization:} The framework employs two-level aggregation with local synchronization every 10 epochs and global synchronization every 50 epochs. This hierarchical structure reduces communication overhead while maintaining convergence stability through differential privacy (noise standard deviation 1.0) and gradient clipping (maximum norm 0.5). The design supports scalable federated learning without requiring fixed cluster assignments.
	
	\item \textbf{Lightweight neural architecture for real-time orchestration:} The PPO agent employs a compact actor-critic network with 16 hidden units per layer, enabling sub-millisecond orchestration decisions while maintaining high reward performance. The architecture uses prioritized experience replay (buffer size 20,000) and entropy regularization (coefficient 0.01) to balance exploration and exploitation.
	
	\item \textbf{Comprehensive experimental validation:} We implement FedSemGNN in EdgeSimPy~\cite{edgesimpy}, a high-fidelity edge computing simulator, with detailed instrumentation of control-plane latency, reward trajectories, semantic fidelity, power consumption, communication footprint, and migration dynamics. The base evaluation uses a heterogeneous edge infrastructure with mobile users following realistic mobility patterns over 1,000 orchestration timesteps.
	
	\item \textbf{Systematic scalability analysis:} Scalability from base topology to 10,000 nodes is validated through mathematical extrapolation using power-law and logarithmic scaling models fitted to empirical measurements. This methodology enables principled large-scale analysis without prohibitive computational overhead. FedSemGNN is compared against four competitive baselines: FlatFedPPO, HierFedPPO, HSQF, and RandomPlacement.
	
	\item \textbf{Superior performance across all metrics:} FedSemGNN achieves normalized reward of 0.91 (8.4 percent higher than FlatFedPPO), orchestration latency of 0.36 ms per decision (315 times faster than FlatFedPPO), communication efficiency of 1.394 reward per MB (174 times higher than FlatFedPPO), 100 percent semantic fidelity with zero migrations (compared to 6,840 in FlatFedPPO), and energy consumption of 72.1 W (36.2 times lower than FlatFedPPO).
\end{itemize}

Unlike previous methods that treat semantics or topology heuristically, FedSemGNN unifies both within a scalable federated RL framework with mathematically grounded reward optimization. The reward function jointly minimizes power consumption and latency while encouraging semantic fidelity through state representation design. By embedding semantic structure into the learning process and encoding topological constraints through graph neural networks, FedSemGNN supports responsive orchestration with minimal latency and communication cost. The system demonstrates real-time operation under 0.5 ms per decision, offering a viable foundation for control-plane design in 6G edge computing.

\noindent\textbf{Paper organization.} The remainder of this paper is structured as follows: Section~\ref{sec:related} reviews related work in federated reinforcement learning, semantic scheduling, and graph neural network based orchestration. Section~\ref{sec:system_architecture} presents the FedSemGNN system architecture including semantic embedding, GNN topology encoding, and hierarchical PPO coordination. Section~\ref{sec:methodology} details the proposed methodology with mathematical formulations for reward optimization, policy learning, and federated aggregation. Section~\ref{sec:results} discusses empirical evaluations across baseline comparisons and scalability analysis. Section~\ref{sec:conclusion} concludes the paper, and Section~\ref{sec:limitations_future} outlines limitations and future research directions.
